\documentclass[dvipdfmx,uplatex,11pt]{beamer}

% --- 日本語設定 ---
\usepackage[utf8]{inputenc}
\usepackage[T1]{fontenc}
\usepackage{bxdpx-beamer} % dvipdfmx + beamer の安定化
\usepackage{pxjahyper}   % 日本語しおり対策
\usepackage{graphicx}
\usepackage{xcolor}

% --- デザイン設定 ---
\usetheme{Metropolis} % シンプルでモダンなテーマ
\usecolortheme{default}
\usefonttheme{professionalfonts}
\setbeamertemplate{navigation symbols}{}
\setbeamertemplate{footline}[frame number]

% --- 数学・記号 ---
\usepackage{amsmath,amsfonts,bm}
\usepackage{mathtools}

\usepackage{tikz}
\usetikzlibrary{arrows.meta,positioning}
\usepackage{pgfplots}
\pgfplotsset{compat=1.18}

% --- 枠付きボックス ---


\usepackage[most]{tcolorbox}

\newtcolorbox{d2lbox}[1]{
  enhanced,
  colback=black!7,
  colframe=black!7,
  colbacktitle=black!22,
  coltitle=black!80,
  title=#1,
  fonttitle=\bfseries,
  boxrule=0pt,
  arc=0pt,
  outer arc=0pt,
  left=2mm,
  right=2mm,
  top=0mm,
  bottom=0mm,
  toptitle=1mm,
  bottomtitle=1mm
}
\newtcolorbox{d2lplainbox}{
  enhanced,
  colback=black!7,
  colframe=black!7,
  boxrule=0pt,
  arc=0pt,
  outer arc=0pt,
  left=2mm,
  right=2mm,
  top=0mm,
  bottom=0mm
}

% --- コード表示 ---
\usepackage{listings}
\lstdefinestyle{pytorch}{
  language=Python,
  basicstyle=\ttfamily\scriptsize,
  keywordstyle=\color{blue!70!black}\bfseries,
  commentstyle=\color{green!40!black},
  stringstyle=\color{orange!60!black},
  showstringspaces=false,
  breaklines=true,
  frame=single,
  framerule=0.2pt,
  rulecolor=\color{black!25},
  backgroundcolor=\color{black!3},
  columns=fullflexible,
  keepspaces=true,
  xleftmargin=2mm,
  xrightmargin=2mm
}
\lstset{style=pytorch}

% --- 強調用 ---
\newcommand{\key}[1]{\alert{#1}}
\newcommand{\shape}[1]{\texttt{#1}}
\newcommand{\code}[1]{\texttt{#1}}


% --- タイトル情報 ---
\title{ソフトマックス回帰：分類のための線形ニューラルネットワーク}
\subtitle{Dive into Deep Learning 輪読ゼミ}
\author{山口 海}
\date{\today}

\begin{document}

\maketitle
%================================================

%================================================
\begin{frame}[allowframebreaks]{目次}
  \tableofcontents[
    sectionstyle=show,
    subsectionstyle=show/show/show
  ]
\end{frame}

%================================================

%================================================
\section{4.1. ソフトマックス回帰の概要}
\subsection{4.1.1. 回帰から分類へ}
\begin{frame}{4.1.1. 回帰から分類へ}
線形回帰では，入力から\alert{連続値}を予測した．

\begin{block}{例：回帰}
家の価格はいくらか？\\
患者は何日入院するか？
\end{block}

ソフトマックス回帰では，入力から\alert{カテゴリ}を予測する．

\begin{block}{例：分類}
この画像は猫・犬・鶏のどれか？\\
このメールは迷惑メールか？
\end{block}
\begin{d2lplainbox}
  分類では，単に数値を出すだけでなく，\\
\alert{各クラスに属する確率}として解釈できる出力がほしい
\end{d2lplainbox}
\end{frame}

%================================================

%================================================
\begin{frame}{分類には2つの見方がある}
\begin{block}{ハードな分類}
入力を1つのクラスに割り当てる．
\[
  \text{image} \longmapsto \text{dog}
\]
\end{block}

\begin{block}{ソフトな分類}
各クラスに属する確率を出す．
\[
  \text{image} \longmapsto
  \begin{cases}
    P(\text{cat}) = 0.1\\
    P(\text{dog}) = 0.8\\
    P(\text{chicken}) = 0.1
  \end{cases}
\]
\end{block}

\vspace{0.5em}
\centering
ソフトマックス回帰では，この\alert{確率的な出力}を作る．
\end{frame}

%================================================


\begin{frame}{単純な画像分類問題}

ここでは，$2 \times 2$ のグレースケール画像を考える．

\[
\begin{pmatrix}
x_1 & x_2\\
x_3 & x_4
\end{pmatrix}
\quad \longmapsto \quad
(x_1,x_2,x_3,x_4)
\]

\begin{itemize}
  \item 各画素値を1つの特徴量として扱う
  \item 入力は4次元ベクトル
  \[
  \bm{x} = (x_1,x_2,x_3,x_4)^\top
  \]
  \item 出力クラスは3つ
  \[
  \{\text{cat},\ \text{chicken},\ \text{dog}\}
  \]
\end{itemize}
各画像は，この3つのうち\alert{どれか1つ}に属すると仮定する．

\end{frame}
%================================================
\begin{frame}{ラベルの表し方}
カテゴリを単に整数で表すこともできる．

\[
  y \in \{1,2,3\}
\]

\vspace{1em}
しかし，この表し方だと問題がある．

\begin{alertblock}{整数ラベルの注意点}
\begin{itemize}
  \item cat < chicken < dog という意図しない大小関係が生まれる
  \item 多くの分類問題では，クラス間に自然な順序はない
\end{itemize}
\end{alertblock}

\begin{block}{そこで}
画像識別や自然言語処理でよく使われるのが，
\alert{one-hot 表現}である．
\end{block}


\end{frame}

%================================================
\begin{frame}{one-hot エンコーディング}
分類では，ラベルをベクトルで表すことが多い．

\[
y \in \{(1,0,0),\ (0,1,0),\ (0,0,1)\}
\]

\vspace{1em}
\begin{center}
\begin{tabular}{c|c}
クラス & one-hot ベクトル \\ \hline
cat & $(1,0,0)$ \\
chicken & $(0,1,0)$ \\
dog & $(0,0,1)$
\end{tabular}
\end{center}

\vspace{1em}
\centering
正解クラスの成分だけが\alert{1}，それ以外は\alert{0}になる．
\end{frame}

%================================================
\subsection{4.1.1.1. 線形モデル}
\begin{frame}{4.1.1.1. 線形モデル：クラスごとにスコアを出す}
入力が4次元，クラスが3種類なので，出力も3つ用意する．

\[
\begin{aligned}
o_1 &= x_1 w_{11} + x_2 w_{12} + x_3 w_{13} + x_4 w_{14} + b_1,\\
o_2 &= x_1 w_{21} + x_2 w_{22} + x_3 w_{23} + x_4 w_{24} + b_2,\\
o_3 &= x_1 w_{31} + x_2 w_{32} + x_3 w_{33} + x_4 w_{34} + b_3.
\end{aligned}
\]

\begin{center}
  \includegraphics[width=0.7\linewidth]{linear_model.png}
\end{center}

\begin{itemize}
  \item $w_{ij}$ は入力 $x_j$ がクラス $i$ に与える重み
  \item $b_i$ はクラス $i$ のバイアス項
  \item $o_i$ はクラス $i$ のスコア
\end{itemize}
\end{frame}

%================================================
\begin{frame}{線形モデルの行列表記}
先ほどの式は，行列を使うと
\[
\begin{pmatrix}
o_1\\
o_2\\
o_3
\end{pmatrix}
=
\begin{pmatrix}
w_{11} & w_{12} & w_{13} & w_{14}\\
w_{21} & w_{22} & w_{23} & w_{24}\\
w_{31} & w_{32} & w_{33} & w_{34}
\end{pmatrix}
\begin{pmatrix}
x_1\\
x_2\\
x_3\\
x_4
\end{pmatrix}
+
\begin{pmatrix}b_1\\
b_2\\
b_3
\end{pmatrix}
\]
と書けるが，$b_i$を行列に組み込んで，
\[
\begin{pmatrix}
o_1\\
o_2\\
o_3
\end{pmatrix}
=
\begin{pmatrix}w_{11} & w_{12} & w_{13} & w_{14} & b_1\\
w_{21} & w_{22} & w_{23} & w_{24} & b_2\\
w_{31} & w_{32} & w_{33} & w_{34} & b_3
\end{pmatrix}
\begin{pmatrix}x_1\\
x_2\\
x_3\\
x_4\\
1
\end{pmatrix}
\]と表記することもある．
\end{frame}

%================================================
\begin{frame}{なぜそのまま確率にできないのか}
線形モデルの出力 $\mathbf{o}$ をそのまま確率とみなしたいが，問題がある．

\begin{alertblock}{確率として必要な条件}
\begin{itemize}
  \item 各成分が非負である：$\hat{y}_i \geq 0$
  \item 全クラスの和が1である：$\sum_i \hat{y}_i = 1$
\end{itemize}
($\hat{y}_i$ はクラス $i$ であると予測する確率)
\end{alertblock}

\vspace{1em}
しかし，線形モデルの出力 $o_i$ は，
\[
  o_i < 0
  \quad \text{や} \quad
  \sum_i o_i \neq 1
\]
になりうる．

\vspace{1em}
\centering
そこで，スコアを確率に変換する関数が必要になる．
\end{frame}

%================================================
\subsection{4.1.1.2. ソフトマックス関数}
\begin{frame}{ソフトマックス関数}
確率ベクトルを $\hat{\mathbf{y}}$ と書くと，

\[
\hat{\mathbf{y}} = \mathrm{softmax}(\mathbf{o})
\]
となり，ソフトマックスは，スコア $\mathbf{o}$ を確率ベクトルに変換する．
ここで，$\mathbf{P(y=i \mid \mathbf{x})} \propto \exp(o_i)$ として,

\[
\hat{y}_i
=
\frac{\exp(o_i)}
{\sum_j \exp(o_j)}
\]

とソフトマックス関数を導入する．

\vspace{1em}
\begin{block}{ソフトマックス関数の特徴}
\begin{itemize}
  \item 指数関数により，すべての成分を正にする
  \item 大きいスコアほど，大きい確率になる
\end{itemize}
\end{block}
\end{frame}

%================================================
\begin{frame}{ソフトマックスのイメージ}
例えば，線形モデルの出力が

\[
\mathbf{o} = (1.2,\ 0.3,\ 2.1)
\]

だったとする．

ソフトマックスを通すと，
\[
\hat{\mathbf{y}}
=
\mathrm{softmax}(\mathbf{o})
=
(0.26,\ 0.11,\ 0.63)
\]

のように，確率として解釈できる値になる．


\begin{alertblock}{予測クラス}
最も確率が大きいクラスを選ぶ．
\[
\hat{y}_3 = 0.63
\quad \Rightarrow \quad
\text{dog}
\]
\end{alertblock}
\end{frame}

%================================================
\begin{frame}{ソフトマックスは最大クラスを変えない}
ソフトマックスはスコアを確率に変換するが，
\alert{大小関係}は変えない．
\[
\operatorname*{argmax}_j \hat{y}_j
=
\operatorname*{argmax}_j o_j
\]

\vspace{1em}
つまり，予測クラスだけを知りたいなら，

\[
\mathbf{o} = (1.2,\ 0.3,\ 2.1)
\]

の時点で，最大の成分 $o_3$ を選べばよい．

\vspace{1em}
\begin{block}{重要}
ソフトマックスは「どのクラスが最大か」を変えずに，\\
スコアを\alert{確率として解釈できる形}に変換する．
\end{block}
\end{frame}

%================================================


%================================================
%================================================
\begin{frame}{ソフトマックスとボルツマン分布}
ソフトマックスは，統計力学に現れる
\alert{ボルツマン分布}とよく似た形をしている．

\vspace{0.5em}

\begin{columns}[T]
\column{0.48\textwidth}
\begin{block}{ソフトマックス}

\[
\hat{y}_i
=
\frac{\exp(o_i)}
{\sum_j \exp(o_j)}
\]
\end{block}

\column{0.48\textwidth}
\begin{block}{ボルツマン分布}

\[
P(i)
=
\frac{\exp\left(-\frac{E_i}{k_{\mathrm{B}}T}\right)}
{\sum_j \exp\left(-\frac{E_j}{k_{\mathrm{B}}T}\right)}
\]
\end{block}
\end{columns}

\vspace{1em}
\centering
どちらも「指数をとって，総和で割る」ことで\\
\alert{確率分布}を作っている．
以下，$\beta = 1/(k_{\mathrm{B}}T)$と置き，$\beta$は逆温度と呼ばれる．
\end{frame}

%================================================
\begin{frame}{ボルツマン分布の意味}
ボルツマン分布では，量子状態 $i$ が出現する確率を$P(i) \propto \exp(-\beta E_i)$ を仮定して，

\[
P(i)
=
\frac{1}{Z} \exp(-\beta E_i)
\]
と表す．ここで，分母は規格化定数で，
\[
Z
=
\sum_j \exp\left(- \beta E_j\right)
\]
であり、分配関数と呼ばれる．



\begin{itemize}
  \item $E_i$：状態 $i$ のエネルギー
  \item $T$：温度
  \item $k_{\mathrm{B}}$：ボルツマン定数
  \item $Z$：分配関数
\end{itemize}




\end{frame}

%================================================
\begin{frame}{ソフトマックスとの対応関係}
ソフトマックスでは，スコア $o_i$ が大きいクラスほど
予測確率が大きくなる．

\begin{columns}[T]
\column{0.48\textwidth}
\begin{block}{ソフトマックス}

\[
\hat{y}_i
=
\frac{\exp(o_i)}
{\sum_j \exp(o_j)}
\]
\end{block}

\column{0.48\textwidth}
\begin{block}{ボルツマン分布}

\[
P(i)
=
\frac{\exp\left(- \beta E_i\right)}
{\sum_j \exp\left(- \beta E_j\right)}
\]
\end{block}
\end{columns}

ボルツマン分布と比較すると，

\[
o_i
\quad \leftrightarrow \quad
-\beta E_i
\]

と対応していると見なせる．
\begin{itemize}
  \item ソフトマックス：スコアが高いほど確率が大きい
  \item ボルツマン分布：エネルギーが低いほど確率が大きい
\end{itemize}
\end{frame}

%================================================
\begin{frame}{分配関数の対応関係}
\begin{columns}
\column{0.48\textwidth}
\begin{block}{ソフトマックスの分配関数}
\[
\sum_j \exp(o_j)
\]
\end{block}
\column{0.48\textwidth}
\begin{block}{ボルツマン分布の分配関数}
\[
Z
=
\sum_j \exp\left(- \beta E_j\right)
\]
\end{block}
\end{columns}

統計力学において，分配関数 $Z$ を求めることができれば，\\
その系の熱力学のすべてがわかると言われる．
\vspace{1em}

    \begin{block}{自由エネルギー $F$}
        \[
            F = -\frac{1}{\beta} \ln Z
        \]
    \end{block}

    \begin{itemize}
        \item 分配関数 $Z$ を求めることができれば，これを微分するだけでエントロピーや内部エネルギーなど，\textbf{その系の熱力学のすべてがわかる}．
    \end{itemize}
\end{frame}


%================================================
\subsection{4.1.1.3. ベクトル化}
\begin{frame}{4.1.1.3. ベクトル化}
計算効率のため，1つのデータごとに計算するのではなく，
$n$個の\alert{ミニバッチ}に対してまとめて計算する．
\[
\mathbf{X} \in \mathbb{R}^{n \times d},
\quad
\mathbf{W} \in \mathbb{R}^{d \times q},
\quad
\mathbf{b} \in \mathbb{R}^{1 \times q}
\]

\[
\begin{aligned}
\mathbf{O}
&=
\mathbf{X}\mathbf{W}
+
\mathbf{b},
\quad
\mathbf{O} \in \mathbb{R}^{n \times q} \\
\hat{\mathbf{Y}}
&=
\mathrm{softmax}(\mathbf{O})
\end{aligned}
\]

\begin{columns}[T]
\column{0.45\textwidth}
\begin{itemize}
  \item $n$：ミニバッチサイズ
  \item $d$：入力特徴量の数
  \item $q$：クラス数
\end{itemize}

\column{0.55\textwidth}
\begin{d2lbox}{計算の流れ}
\begin{enumerate}
  \item $\mathbf{X}$ をまとめて線形変換
  \item クラスごとのスコア $\mathbf{O}$ を出す
  \item 各行に softmax をかける
\end{enumerate}
\end{d2lbox}
\end{columns}


\end{frame}

%================================================
\begin{frame}{ミニバッチを行列で表す}
入力特徴量の数を $d$，ミニバッチサイズを $n$ とする．
$n$ 個の入力データを縦に並べる．
\[
\mathbf{X}
=
\begin{pmatrix}
- & (\mathbf{x}^{(1)})^\top & -\\
  & \vdots & \\
- & (\mathbf{x}^{(n)})^\top & -
\end{pmatrix}
\in \mathbb{R}^{n \times d}
\]
$\mathbf{X}$ の各行が，1つの入力データを表す．
重み行列$\mathbf{W} \in \mathbb{R}^{d \times q}$を掛け，１つのデータに対するクラスごとの補正であるバイアス $\mathbf{b} \in \mathbb{R}^{1 \times q}$ を足すと，
\[
\mathbf{O}
=
\begin{pmatrix}
- & (\mathbf{x}^{(1)})^\top \mathbf{W} & -\\
  & \vdots & \\
- & (\mathbf{x}^{(n)})^\top \mathbf{W} & -
\end{pmatrix}
+
\begin{pmatrix}
- & \mathbf{b} & -\\
  & \vdots & \\
- & \mathbf{b} & -
\end{pmatrix}
\]
スコア行列 $\mathbf{O} \in \mathbb{R}^{n \times q}$ を得る．
このような自動的な拡張を\alert{ブロードキャスト}という．

\end{frame}

%================================================

\begin{frame}{ベクトル化されたソフトマックス}
線形モデルで得たスコア行列に対して，softmax をかける．
\[
\hat{\mathbf{Y}}
=
\mathrm{softmax}(\mathbf{O})
\]
\[
\hat{Y}_{ij}
=
\frac{\exp(O_{ij})}
{\sum_{k=1}^{q}\exp(O_{ik})}
\]
softmax は，\alert{各行ごと}に計算される．
\[
\hat{\mathbf{Y}}
=
\begin{pmatrix}
\hat{y}^{(1)}_1 & \hat{y}^{(1)}_2 & \cdots & \hat{y}^{(1)}_q \\
\hat{y}^{(2)}_1 & \hat{y}^{(2)}_2 & \cdots & \hat{y}^{(2)}_q \\
\vdots & \vdots & \ddots & \vdots \\
\hat{y}^{(n)}_1 & \hat{y}^{(n)}_2 & \cdots & \hat{y}^{(n)}_q
\end{pmatrix}
\in \mathbb{R}^{n \times q}
\]
$\hat{y}^{(i)}_j$ は，$i$ 番目のデータがクラス $j$ であると予測する確率を表し，各行の和は1になる．



\end{frame}



% =========================
% 4.1.2 損失関数
% =========================

\subsection{4.1.2. 損失関数}
\begin{frame}{4.1.2. 損失関数}
ソフトマックス回帰では, モデルの出力を
\alert{各クラスに属する確率}として解釈する.

\[
\hat{y}_j = \frac{\exp(o_j)}{\sum_{k=1}^{q} \exp(o_k)}
\]

\[
\hat{\bm y}
=
(\hat y_1,\hat y_2,\dots,\hat y_q),
\qquad
\hat y_j
=
P(y=j\mid \bm x)
\]

\begin{d2lplainbox}
損失関数は, \alert{正解クラスに高い確率を与えるほど小さくなる}ように作りたい.
\end{d2lplainbox}

\end{frame}

%================================================

\begin{frame}{4.1.2. 損失関数}
\[
\hat{\bm y}=(\hat y_{\mathrm{cat}},\hat y_{\mathrm{dog}},\hat y_{\mathrm{bird}}),
\qquad
\bm y=(0,1,0)
\]

ここでは, 正解クラスは\alert{犬}である.

\begin{columns}[T]
\begin{column}{0.48\textwidth}
\begin{d2lbox}{良い予測}
\[
\hat{\bm y}=(0.1,0.7,0.2)
\]
\[
P(y=\mathrm{dog}\mid \bm x)=0.7
\]
正解クラスである犬に高い確率を与えている.
\[
\Rightarrow \text{損失は小さい}
\]
\end{d2lbox}
\end{column}

\begin{column}{0.48\textwidth}
\begin{d2lbox}{悪い予測}
\[
\hat{\bm y}=(0.8,0.1,0.1)
\]
\[
P(y=\mathrm{dog}\mid \bm x)=0.1
\]
正解クラスである犬に低い確率しか与えていない.
\[
\Rightarrow \text{損失は大きい}
\]
\end{d2lbox}
\end{column}
\end{columns}


\end{frame}


\begin{frame}{最尤推定から損失関数を作る}
データセットを
\[
\{(\bm x^{(1)},y^{(1)}),\dots,(\bm x^{(n)},y^{(n)})\}
\]
とする.

各データが独立にサンプリングされると仮定すると,
データセット全体の条件付き確率は

\[
P(\bm Y\mid \bm X)
=
\prod_{i=1}^{n}
P(y^{(i)}\mid \bm x^{(i)})
\]

と分解できる.

\begin{d2lplainbox}
この量を大きくすることは,
観測された正解ラベルが最も起こりやすくなるように
モデルを調整することに対応する.
\end{d2lplainbox}
\end{frame}

\subsection{4.1.2.1. 対数尤度}
\begin{frame}{負の対数尤度}
尤度は積の形なので, そのままでは扱いにくい.
そこで対数を取る.
\[
\log P(\bm Y\mid \bm X)
=
\log
\prod_{i=1}^{n}
P(y^{(i)}\mid \bm x^{(i)})
=
\sum_{i=1}^{n}
\log P(y^{(i)}\mid \bm x^{(i)})
\]
尤度最大化は, 負の対数尤度最小化と同値である.
\[
-\log P(\bm Y\mid \bm X)
=
\sum_{i=1}^{n}
-\log P(y^{(i)}\mid \bm x^{(i)})
\]
ここで$i$番目のデータ1つに対する損失を
\[
l(\bm y^{(i)},\hat{\bm y}^{(i)})
=
-\log P(y^{(i)}\mid \bm x^{(i)})
\]
と定義する．
\end{frame}


\begin{frame}{one-hot ラベルと交差エントロピー}
正解が第 $k$ クラスであるとすると, one-hot ベクトルは
\[\bm y=(0,\dots,0,1,0,\dots,0)
\]
となる. このとき
\[
  - \sum_{j=1}^{q} y_j \log \hat{y}_j
  =
  - \log \hat{y}_k
\]
となるので，
\[
\boxed{
l(\bm y,\hat{\bm y})
=
-\sum_{j=1}^{q}
y_j\log \hat y_j
}
\]
となる．この損失関数は\alert{交差エントロピー損失}と呼ばれる．

\end{frame}

%================================================

\begin{frame}[t]{具体例: 正解クラスの確率と交差エントロピー損失}
正解ラベルが第2クラスであるとする．
\[
\bm y=(0,1,0)
\]
\vspace{-8mm}
\begin{columns}[T,onlytextwidth]
\begin{column}{0.55\textwidth}
\begin{d2lbox}{良い予測}
\[
\hat{\bm y}=(0.1,0.7,0.2)
\]
正解クラスの確率：$\hat y_2=0.7$ \\
このときの損失は
\[
\begin{split}
&l(\bm{y},\hat{\bm{y}})
= - \sum_{j=1}^{3} y_j \log \hat{y}_j = - (0\log 0.1 \\
&+ 1\log 0.7 + 0\log 0.2) \approx 0.357
\end{split}
\]
\[
\Rightarrow \text{損失は小さい}
\]
\end{d2lbox}
\end{column}

\begin{column}{0.4\textwidth}
\begin{d2lbox}{悪い予測}
\[
\hat{\bm y}=(0.8,0.1,0.1)
\]
\[
\hat y_2=0.1
\]
損失は
\[
l(\bm y,\hat{\bm y})
=
-\log 0.1
\approx 2.303
\]
\[
\Rightarrow \text{損失は大きい}
\]
\end{d2lbox}
\end{column}
\end{columns}

\end{frame}

%================================================
\subsection{4.1.2.2. ソフトマックスと交差エントロピー損失}
\begin{frame}{4.1.2.2. ソフトマックスと交差エントロピー損失}
ソフトマックス回帰では, ロジット $\bm o$ を確率に変換する.

\[
\hat y_j
=
\operatorname{softmax}(\bm o)_j
=
\frac{\exp(o_j)}
{\sum_{k=1}^{q}\exp(o_k)}
\]

交差エントロピー損失は

\[
l(\bm y,\hat{\bm y})
=
-\sum_{j=1}^{q}y_j\log \hat y_j
\]

であった.

\begin{d2lplainbox}
ここでは, この2つを組み合わせると何が起こるかを見る.
特に, 勾配が非常に簡単な形になる.
\end{d2lplainbox}
\end{frame}


\begin{frame}{交差エントロピーにソフトマックスを代入する}
交差エントロピー損失：$
l(\bm y,\hat{\bm y})
=
-\sum_{j=1}^{q}y_j\log \hat y_j
$
に

$
\hat y_j
=
\frac{\exp(o_j)}
{\sum_{k=1}^{q}\exp(o_k)}
$
を代入して$\sum_{j=1}^{q}y_j=1$を使うと,
{\small
\[
\begin{aligned}
  l(\bm y,\hat{\bm y})
&=
-\sum_{j=1}^{q}
y_j
\log
\frac{\exp(o_j)}
{\sum_{k=1}^{q}\exp(o_k)}
= 
-\sum_{j=1}^{q}
y_j
\left(
o_j
-
\log\sum_{k=1}^{q}\exp(o_k)
\right) \\
&= \sum_{j=1}^{q}
y_j
\log\sum_{k=1}^{q}\exp(o_k)
-
\sum_{j=1}^{q}y_j o_j
=
\log\sum_{k=1}^{q}\exp(o_k)
-
\sum_{j=1}^{q}y_j o_j 
\end{aligned}
\]}
となる．
任意のロジット$o_j$に対して勾配を取る．
{\small
\[
\begin{aligned}
  \frac{\partial}{\partial o_j} l(\bm y,\hat{\bm y})
&=
\frac{\partial}{\partial o_j}
\log\sum_{k=1}^{q}\exp(o_k)
-
\frac{\partial}{\partial o_j}\sum_{k=1}^{q}y_k o_k \\
&= \frac{\exp(o_j)}
{\sum_{k=1}^{q}\exp(o_k)}
-
y_j 
= \mathrm{softmax}(\bm o)_j - y_j \\
&= \hat y_j - y_j
\end{aligned}
\]}

\end{frame}
%================================================


\begin{frame}{勾配の意味}
\[
\boxed{
\frac{\partial l}{\partial o_j}
=
\hat y_j-y_j
}
\]
つまり, 勾配は「$\text{予測確率} - \text{正解ラベル}$」である. \\
例：$\bm y = (0,1,0)$ 
\small{
\begin{columns}[T,onlytextwidth]
\begin{column}{0.48\textwidth}
\begin{d2lbox}{良い予測}
\[
\hat{\bm y}=(0.1,0.7,0.2)
\]
\[\hat{\bm y}-\bm y
=
(0.1,-0.3,0.2)
\]
正解クラスの修正は小さい.
\end{d2lbox}
\end{column}

\begin{column}{0.48\textwidth}
\begin{d2lbox}{悪い予測}
\[
\hat{\bm y}=(0.8,0.1,0.1)
\]
\[\hat{\bm y}-\bm y
=
(0.8,-0.9,0.1)
\]
 強く修正される.
\end{d2lbox}
\end{column}
\end{columns}}
\[
\text{勾配降下法：}\quad \bm{o} \leftarrow \bm{o} - \eta \frac{\partial l}{\partial \bm{o}}
\]
\begin{d2lplainbox}
  正解クラスのスコアを上げ，不正解のスコアを下げるように更新
\end{d2lplainbox}

\end{frame}
%================================================

\subsection{4.1.3. 情報理論の基礎}
\begin{frame}{情報量: 驚きの大きさ}
情報理論では, 事象が起きたときの\alert{驚きの大きさ}を
情報量として考える.
\[
\text{情報量}
=
\text{驚きの大きさ}
\]
事象 $j$ が起こる確率を $P(j)$ とすると, その情報量を
\[
\boxed{
I(j)
\coloneqq
-\log P(j)
}
\]
と定義する.
\begin{d2lplainbox}
まれな事象は貴重なサンプルとなり情報量は大きい．
\[
\text{確率が低い} \Leftrightarrow \text{情報量が大きい}
\]
\end{d2lplainbox}
\end{frame}
%================================================

\subsection{4.1.3.1. 情報エントロピー}
\begin{frame}{情報エントロピー: 情報量の期待値}
分布 $P$ に従って事象 $j$ が起こるとする.
各事象の情報量は
\[
I(j)=-\log P(j)
\]
である.
したがって, 平均的な情報量は期待値として
\[
\mathbb{E}_{j\sim P}[I(j)]
=
\sum_j P(j)I(j)
\]
と書ける.これを\alert{情報エントロピー}という．
\[
\boxed{
\begin{aligned}
H[P]
&=
\sum_j P(j)(-\log P(j)) \\
&=
-\sum_j P(j)\log P(j)
\end{aligned}
}
\]

\end{frame}

%================================================
\subsection{4.1.3.3. 交差エントロピーの再考}
\begin{frame}{4.1.3.3. 交差エントロピーの再考}
真の分布を $P$, モデルの予測分布を $Q$ として交差エントロピーを定義する．
\[
H(P,Q)
\coloneqq
-\sum_j P(j)\log Q(j)
\]
分類の交差エントロピー損失は$
P(j)=y_j,
\quad
Q(j)=\hat y_j
$
で対応．
\[
H(P,Q)
=
-\sum_{j=1}^{q}y_j\log\hat y_j
=
l(\bm y,\hat{\bm y})
\]
\begin{d2lplainbox}
交差エントロピー損失は,
真の分布$P$に従って起こる事象を，予測分布$Q$で説明するときの
平均的情報量である．
\end{d2lplainbox}
予測分布 $Q$ が真の分布 $P$ に近いほど小さくなる.つまり$Q=P$
のとき, 交差エントロピーは最小になる．
\[
H(P,Q)=H(P,P)=H(P)
\]
\end{frame}


%================================================
\section{4.2. 画像分類データセットと実践}
\subsection{4.2.1. データセットの読み込み}
\begin{frame}{4.2. 画像分類データセット}
この節では，画像分類の実験に用いる
\alert{Fashion-MNIST} を準備する．

\begin{d2lbox}{この節の位置づけ}
\begin{itemize}
  \item モデルの理論ではなく，データの準備を扱う
  \item 画像とラベルの組を読み込む
  \item ミニバッチとしてモデルに入力できる形にする
\end{itemize}
\end{d2lbox}

以降のソフトマックス回帰では，このデータを使って分類を行う．

\centering
\large
データセット $\longrightarrow$ ミニバッチ $\longrightarrow$ 分類モデル
\end{frame}

\begin{frame}{Fashion-MNISTとは}
\key{Fashion-MNIST}は，衣類画像を10クラスに分類するデータセットである．

\vspace{-1.5em}
\begin{columns}[T,totalwidth=\textwidth]
\begin{column}{0.52\textwidth}
\begin{d2lbox}{特徴}
\begin{itemize}
  \item 学習データ：$60000$ 枚
  \item テストデータ：$10000$ 枚
  \item クラス数：$10$
  \item 画像サイズ：$28 \times 28$
  \item 白黒画像なのでチャネル数は $1$
\end{itemize}
\end{d2lbox}
\end{column}
\begin{column}{0.44\textwidth}
\begin{d2lplainbox}
\small
\begin{center}
\begin{tabular}{c|c}
ラベル & クラス名 \\ \hline
0 & t-shirt \\
1 & trouser \\
2 & pullover \\
3 & dress \\
4 & coat \\
5 & sandal \\
6 & shirt \\
7 & sneaker \\
8 & bag \\
9 & ankle boot
\end{tabular}
\end{center}
\end{d2lplainbox}
\end{column}
\end{columns}
\[
\text{画像 } \bm{x}
\quad \longrightarrow \quad
\text{ラベル } y \in \{0,1,\dots,9\}
\]
\end{frame}

\begin{frame}[fragile]{PyTorchでデータセットを読み込む}
PyTorch では，\texttt{torchvision.datasets} を用いて
Fashion-MNIST を読み込める．
\begin{lstlisting}
class FashionMNIST(d2l.DataModule):  #@save
    """The Fashion-MNIST dataset."""
    def __init__(self, batch_size=64, resize=(28, 28)):
        super().__init__()
        self.save_hyperparameters()
        trans = transforms.Compose([transforms.Resize(resize),
                                    transforms.ToTensor()])
        self.train = torchvision.datasets.FashionMNIST(
            root=self.root, train=True, transform=trans, download=True)
        self.val = torchvision.datasets.FashionMNIST(
            root=self.root, train=False, transform=trans, download=True)

\end{lstlisting}
\begin{itemize}
  \item \texttt{train=True}: 学習用データを使う
  \item \texttt{train=False}: 評価用データ,検証用データを使う
  \item \texttt{ToTensor()}: 画像を数値テンソルに変換
\end{itemize}
\end{frame}

\begin{frame}[fragile]{テンソルの形の意味}
1枚の画像の形は，
\begin{lstlisting}
data.train[0][0].shape

# torch.Size([1, 32, 32])
\end{lstlisting}

\begin{d2lbox}{画像テンソルの形}
\[
\text{image} \in \mathbb{R}^{c \times h \times w}
\]
\[
\shape{[1, 32, 32]}
=
\texttt{[channel 1, height 32, width 32]}
\]
$32 \times 32$ の白黒画像なので，チャネル数は $1$ である．
\end{d2lbox}
\begin{alertblock}{注意}
元画像は$28\times 28$だが，ここではD2Lに合わせて$32\times32$にリサイズしている．
\end{alertblock}

\end{frame}

\subsection{4.2.2 ミニバッチの読み込み}
\begin{frame}[fragile]{4.2.2 ミニバッチの読み込み}
学習では，複数枚をまとめた\key{ミニバッチ}で処理する．
\begin{lstlisting}
@d2l.add_to_class(FashionMNIST)  #@save
def get_dataloader(self, train):
    data = self.train if train else self.val
    return torch.utils.data.DataLoader(data, self.batch_size, shuffle=train,
                                       num_workers=self.num_workers)
X, y = next(iter(data.train_dataloader()))
print(X.shape, X.dtype, y.shape, y.dtype)

# torch.Size([64, 1, 32, 32]) torch.float32 torch.Size([64]) torch.int64
\end{lstlisting}
\vspace{-1em}
\begin{d2lplainbox}
  ミニバッチの形状：$n=64$, $c=1$, $h=w=32$
\[
\bm{X} \in \mathbb{R}^{n \times c \times h \times w}
\]
テンソルの形：$\texttt{[batch size, channel, height, width]}$
$32 \times 32$ の白黒画像を64枚まとめている．
\end{d2lplainbox}
\end{frame}

\subsection{4.2.3 画像の可視化}
\begin{frame}[fragile]{4.2.3 画像の可視化}
データが正しく読み込めているかを確認するために，
画像を表示する．
\begin{lstlisting}
@d2l.add_to_class(FashionMNIST)
def visualize(self, batch, nrows=1, ncols=8, labels=[]):
    X, y = batch
    if not labels:
        labels = self.text_labels(y)
    d2l.show_images(X.squeeze(1), nrows, ncols,
                    titles=labels)
\end{lstlisting}

画像とラベルが対応しているか，リサイズやテンソル変換が意図通りかを確認する．
\end{frame}

\subsection{4.2.4. まとめ}
\begin{frame}{4.2節のまとめ}
\begin{d2lbox}{この節の要点}
\begin{itemize}
  \item Fashion-MNISTは10クラスの衣類画像分類データセットである．
  \item PyTorchでは\texttt{torchvision.datasets.FashionMNIST}で読み込める．
  \item 画像は\shape{[チャネル数, 高さ, 幅]}のテンソルとして扱う．
  \item ミニバッチでは\shape{[バッチサイズ, チャネル数, 高さ, 幅]}になる．
  \item \texttt{DataLoader}は，シャッフル・ミニバッチ化・並列読み込みを担う．
\end{itemize}
\end{d2lbox}
\end{frame}

\section{4.3. 分類の基礎モデル}

\begin{frame}{4.3. 分類の基礎モデル}
\begin{d2lbox}{この節でやること}
分類問題で何度も使う処理を，共通の部品としてまとめる．
\end{d2lbox}

\vspace{3mm}

分類モデルでは，毎回ほぼ同じ流れが現れる：
\[
\text{入力 } \bm X
\longrightarrow
\text{モデル}
\longrightarrow
\text{予測 } \hat{\bm Y}
\longrightarrow
\text{損失・精度の計算}
\]

\vspace{3mm}

\begin{itemize}
  \item \alert{損失}：学習のために最小化する量
  \item \alert{精度}：分類器がどれくらい正解したかを表す指標
\end{itemize}

\vspace{2mm}

\[
\mathrm{accuracy}
=
\frac{\text{正解したデータ数}}{\text{全データ数}}
\]
\end{frame}

\subsection{4.3.1. Classifier クラス}
\begin{frame}[fragile]{4.3.1. \texttt{Classifier} クラス}
分類モデル用の基底クラスを用意する．

\vspace{2mm}

\begin{lstlisting}
class Classifier(d2l.Module):  #@save
    """分類モデルの基底クラス。"""
    def validation_step(self, batch):
        Y_hat = self(*batch[:-1])
        self.plot('loss', self.loss(Y_hat, batch[-1]), train=False)
        self.plot('acc', self.accuracy(Y_hat, batch[-1]), train=False)
\end{lstlisting}
\begin{d2lplainbox}
検証データに対して，
\[
\text{予測}
\rightarrow
\text{損失の計算}
\rightarrow
\text{精度の計算}
\]
をまとめて行う．
\end{d2lplainbox}
\end{frame}

\subsection{4.3.2. 精度}
\begin{frame}{4.3.2. 精度 accuracy}
分類では，予測確率またはスコアから，最終的に1つのクラスを選ぶ．

\begin{center}
\begin{tikzpicture}[node distance=1.6cm,>=Latex]
  \node[draw,rounded corners,fill=blue!6,minimum width=2.5cm,minimum height=0.8cm] (score) {予測スコア $\hat{\bm y}$};
  \node[draw,rounded corners,fill=green!6,minimum width=2.5cm,minimum height=0.8cm,right=of score] (pred) {予測クラス $\hat y$};
  \node[draw,rounded corners,fill=orange!8,minimum width=2.5cm,minimum height=0.8cm,right=of pred] (acc) {正解と比較};
  \draw[->,thick] (score) -- node[above]{\small $\arg\max$} (pred);
  \draw[->,thick] (pred) -- node[above]{\small $\hat y = y$ ?} (acc);
\end{tikzpicture}
\end{center}
\begin{d2lbox}{分類精度}
\[
\mathrm{accuracy}
=
\frac{\text{正しく分類できたデータ数}}{\text{全データ数}}
\]
\end{d2lbox}
これを数式で書くと，
\[
\mathrm{accuracy}
=
\frac{1}{n}
\sum_{i=1}^{n}
\bm{1}(\hat y_i = y_i)
\]
ここで
\vspace{-0.5em}
{\small
\[
\bm{1}(\hat y_i = y_i)
=
\begin{cases}
1 \quad ( \hat y_i = y_i) \\
0  \quad ( \hat y_i \neq y_i)
\end{cases}
\]
}

\end{frame}

\begin{frame}[fragile]{精度の計算：D2Lの実装}
各データについて，最大スコアのクラスを予測クラスとする．

\begin{lstlisting}
@d2l.add_to_class(Classifier)  #@save
def accuracy(self, Y_hat, Y, averaged=True):
    """正しい予測の数を計算する。"""
    Y_hat = d2l.reshape(Y_hat, (-1, Y_hat.shape[-1]))
    preds = d2l.astype(d2l.argmax(Y_hat, axis=1), Y.dtype)
    compare = d2l.astype(preds == d2l.reshape(Y, -1), d2l.float32)
    return d2l.reduce_mean(compare) if averaged else compare
\end{lstlisting}

\begin{d2lbox}{流れ}
\[
\hat{Y} \in \mathbb{R}^{n\times q}
\quad\xrightarrow{\arg\max}\quad
\hat{y}_1,\ldots,\hat{y}_n
\quad\xrightarrow{\hat{y}_i=y_i}\quad
0/1
\]
\end{d2lbox}

\begin{itemize}
  \item \code{averaged=True}: 正解率の平均を返す
  \item \code{averaged=False}: 各データの正誤を0/1で返す
\end{itemize}
\end{frame}

\subsection{4.3.3 まとめ}
\begin{frame}{まとめ}
\begin{d2lbox}{4.3 の要点}
分類モデルでは，予測結果を評価するために
\alert{精度 accuracy} を計算する．
\end{d2lbox}
\begin{itemize}
  \item \texttt{Classifier} クラスは分類モデルの共通処理をまとめる
  \item 検証時には，損失 \texttt{loss} と精度 \texttt{acc} を記録する
  \item 予測クラスは
  \[
  \operatorname*{argmax}_{j}\hat y_j
  \]
  で決める
  \item 精度は
  \[
  \frac{\text{正解数}}{\text{全データ数}}
  \]
  である
\end{itemize}
次節では，この仕組みを使ってソフトマックス回帰を実装する．
\end{frame}

\end{document}